\chapter{Αλγεβρικό Υπόβαθρο: Πεπερασμένα Σώματα και Αριθμητική σε \texorpdfstring{$\GFpoly$}{GF(2)[x]}}
\label{app:finite-fields}

Το παρόν παράρτημα συγκεντρώνει το αλγεβρικό υπόβαθρο των πεπερασμένων σωμάτων και της αριθμητικής πολυωνύμων πάνω στο οποίο εδράζεται το \HQC. Το κυρίως κείμενο (\cref{ch:theoretical-background}) προϋποθέτει αυτό το υλικό, αλλά εστιάζει στον αλγόριθμο και στις βελτιστοποιήσεις του.

\section{Συμβολισμός Πεπερασμένων Σωμάτων}
\label{sec:appendix-field-notation}

Το $\GF{2}$ συμβολίζει το \emph{πεπερασμένο σώμα} (finite field· ιστορικά αποκαλείται και Galois field, εξ ου και ο συμβολισμός GF) δύο στοιχείων $\{0,1\}$.
Ως πεπερασμένο σώμα ορίζεται ένα σύνολο με πεπερασμένο πλήθος στοιχείων, εφοδιασμένο με πρόσθεση και πολλαπλασιασμό που υπακούουν στους συνήθεις κανόνες της αριθμητικής (αντιμεταθετικότητα, προσεταιριστικότητα, επιμεριστικότητα) και στο οποίο κάθε στοιχείο διαθέτει αντίθετο ως προς την πρόσθεση και κάθε μη μηδενικό στοιχείο αντίστροφο ως προς τον πολλαπλασιασμό.
Στο $\GF{2}$ τα δύο μόνα στοιχεία είναι το $0$ και το $1$, η πρόσθεση ορίζεται ως XOR ($1{+}1{=}0$) και ο πολλαπλασιασμός ως AND· έτσι η αριθμητική του σώματος ταυτίζεται με στοιχειώδεις δυαδικές λογικές πράξεις.
Ο $\GF{2}^n$ είναι ο διανυσματικός χώρος των δυαδικών ακολουθιών μήκους $n$ υπεράνω του $\GF{2}$.
Το ότι ο $\GF{2}^n$ ορίζεται \emph{υπεράνω του} $\GF{2}$ σημαίνει ότι οι συντεταγμένες κάθε διανύσματος ανήκουν στο $\GF{2}$ (είναι δηλαδή bits) και ότι τόσο η πρόσθεση διανυσμάτων όσο και ο βαθμωτός πολλαπλασιασμός εκτελούνται κατά συντεταγμένη υπεράνω του $\GF{2}$· πρακτικά, κάθε στοιχείο του $\GF{2}^n$ είναι μια δυαδική ακολουθία $n$ bits και η πρόσθεση δύο τέτοιων διανυσμάτων ισοδυναμεί με bitwise XOR.
Με $\GFpoly$ συμβολίζεται ο δακτύλιος των πολυωνύμων με συντελεστές στο $\GF{2}$, δηλαδή το σύνολο όλων των πολυωνύμων μίας μεταβλητής $x$ με συντελεστές $0$ ή $1$, εφοδιασμένο με τη συνήθη πρόσθεση και πολλαπλασιασμό πολυωνύμων· ο όρος \emph{δακτύλιος} (ring) υποδηλώνει ότι η δομή υποστηρίζει πρόσθεση και πολλαπλασιασμό με τις αναμενόμενες ιδιότητες, χωρίς ωστόσο να εγγυάται την ύπαρξη πολλαπλασιαστικού αντιστρόφου για κάθε στοιχείο.
Με $\GFquot{N}$ συμβολίζεται ο αντίστοιχος \defterm{δακτύλιος-πηλίκο} (quotient ring) modulo $x^N-1$: τα στοιχεία του είναι πολυώνυμα βαθμού $<N$ και, επειδή $x^N \equiv 1$, ο πολλαπλασιασμός ισοδυναμεί με κυκλική συνέλιξη των διανυσμάτων συντελεστών, καθώς οι όροι βαθμού $\geq N$ αναδιπλώνονται στις χαμηλές θέσεις.

\section{Αριθμητική σε \texorpdfstring{$\GFpoly$}{GF(2)[x]}}
\label{sec:gf2x-arithmetic}

Στον δακτύλιο $\GFpoly$, η πρόσθεση ορίζεται ως
XOR κατά συντελεστή και ο πολλαπλασιασμός ως \defterm{carry-less πολλαπλασιασμός}:
πολλαπλασιασμός πολυωνύμων χωρίς κρατούμενο (carry) στις ενδιάμεσες προσθέσεις.
Τυπικά, για δύο πολυώνυμα $a(x) = \sum_{i=0}^{m} a_i x^i$ και
$b(x) = \sum_{j=0}^{n} b_j x^j$ με $a_i, b_j \in \{0, 1\}$, το γινόμενο
ορίζεται ως:
\begin{equation}
  c(x) = a(x) \cdot b(x) = \sum_{k=0}^{m+n} c_k x^k, \quad
  c_k = \bigoplus_{i+j=k} a_i \cdot b_j.
  \label{eq:gf2x-product}
\end{equation}
Η πρόσθεση $\oplus$ αντιστοιχεί σε XOR (πρόσθεση modulo 2), ενώ ο
πολλαπλασιασμός $a_i \cdot b_j$ αντιστοιχεί σε AND.
Η απουσία κρατουμένου διαφοροποιεί ουσιαστικά τον carry-less πολλαπλασιασμό
από τον ακέραιο πολλαπλασιασμό: δεν υπάρχει αλυσιδωτή διάδοση κρατουμένου μεταξύ θέσεων.
